\section*{Chapitre \ref{ch:m1:the-buy-side} --- Le buy side}

\begin{solution}{exo:m1:the-buy-side:1}
Rendements actifs $+1.1$, $-0.6$, $+0.8$, $+1.5$~; moyenne $0.70\%$~; écart-type
empirique $0.91\%$~; $\mathrm{IR} = 0.77$ --- sur quatre observations, ce que
la \cref{prop:m1:the-buy-side:years} juge encore sans valeur.
\end{solution}

\begin{solution}{exo:m1:the-buy-side:2}
$4/0.4^2 = 25$ ans~; $4/1.5^2 = 1.8$ ans~; $4/36$ d'une année, 28 jours
de bourse.
\end{solution}

\begin{solution}{exo:m1:the-buy-side:3}
Fonds indiciel~: $50 \times 10^9 \times 0.0004 = \$20$ millions. \omterm{def:m1:what-a-trading-firm-does:hedge-fund}{Hedge fund}~:
\omterm{def:m1:the-buy-side:fees}{frais de gestion} 40~millions de dollars, \omterm{def:m1:the-buy-side:fees}{commission de performance} $0.2 \times (200 - 40) = \$32$ millions,
soit 72~millions de dollars au total, sur un vingt-cinquième des actifs.
\end{solution}

\begin{solution}{exo:m1:the-buy-side:4}
Les positions supérieures à 5~\% totalisent $9.5 + 18 + 11 = 38.5\% \le 40\%$ et aucune ne dépasse
10~\%~: conforme. Après la hausse, le titre pèse $12.35/102.85 = 12.0\%$~:
au-dessus de 10~\%. Le gérant doit vendre environ un sixième de la position, quel que soit son
avis sur la société --- un ordre sans aucune information.
\end{solution}

\begin{solution}{exo:m1:the-buy-side:5}
$\mathrm{TE} = \sqrt{100 \times (0.005\times0.30)^2} = 1.5\%$. Pour 3~\%, il faut doubler
les poids actifs~: 1.0~\% par position.
\end{solution}

\begin{solution}{exo:m1:the-buy-side:6}
Année 1~: \omterm{def:m1:the-buy-side:fees}{frais de gestion} 2.00, $N' = 128.00$, \omterm{def:m1:the-buy-side:fees}{commission de performance} $0.2\times28 = 5.60$,
clôture 122.40 (pour 100). Année 2~: \omterm{def:m1:the-buy-side:fees}{frais de gestion} 2.45, $N' = 95.47$, sous le
\omterm{def:m1:the-buy-side:fees}{high-water mark} de 122.40~: pas de \omterm{def:m1:the-buy-side:fees}{commission de performance}, clôture 95.47. Année 3~: \omterm{def:m1:the-buy-side:fees}{frais de gestion} 1.91, $N' =
103.11$, toujours sous 122.40~: pas de commission, clôture 103.11.
\end{solution}

\begin{solution}{exo:m1:the-buy-side:7}
«~2 et 20~»~: frais de 35.9 pour 100 investis, valeur finale 145.2. «~1 et 30~»~:
frais 33.2, valeur finale 149.8. L'investisseur qui croit au gérant
préfère «~1 et 30~»~: il paie moins au total ici et paie surtout quand il a été
rendu plus riche. Le gérant peut préférer «~2 et 20~» parce que les deux tiers
de ses revenus (23.4 sur 35.9) ne dépendent alors en rien de la performance, et
paient les salaires durant les sept années passées sous le \omterm{def:m1:the-buy-side:fees}{high-water mark}.
\end{solution}

\begin{solution}{exo:m1:the-buy-side:8}
Sans aucune compétence, le rendement actif moyen sur cinq ans suit $N(0, 4\%/\sqrt5) = N(0,
1.79\%)$. Le vingtième des 400 est le 95\up{e}~centile, $1.645 \times 1.79 =
2.9\%$ par an, un \omterm{def:m1:the-buy-side:ir}{ratio d'information} apparent de $2.9/4 = 0.74$~; les meilleurs
une poignée dépasse 1. Le défaut tient à la sélection~: le classement garantit des
bilans impressionnants parmi du pur bruit, et ces bilans n'ont aucun rapport avec les cinq
années suivantes. Le test honnête est la performance des fonds sélectionnés
\emph{après} la sélection.
\end{solution}

\begin{solution}{pb:m1:the-buy-side:1}
Pour 100~dollars investis.
\textbf{1.} \omterm{def:m1:the-buy-side:fees}{Frais de gestion} 2.00~; $N' = 118.00$~; \omterm{def:m1:the-buy-side:fees}{commission de performance} $0.2 \times 18 = 3.60$~;
clôture 114.40~; \omterm{def:m1:the-buy-side:fees}{high-water mark} 114.40.
\textbf{2.} \omterm{def:m1:the-buy-side:fees}{Frais de gestion} 2.29~; $N' = 114.40\times0.85 - 2.29 = 94.95$~; sous
le \omterm{def:m1:the-buy-side:fees}{high-water mark}, donc pas de \omterm{def:m1:the-buy-side:fees}{commission de performance}~; le \omterm{def:m1:the-buy-side:fees}{high-water mark} reste à 114.40.
\textbf{3.} \omterm{def:m1:the-buy-side:fees}{Frais de gestion} 1.90~; $N' = 102.55 < 114.40$~: pas de commission.
\textbf{4.} L'investisseur affiche $+2.55$~; le gérant a perçu $9.79$.
\textbf{5.} Années 1, 4 et 7.
\textbf{6.} \omterm{def:m1:the-buy-side:fees}{Frais de gestion} 23.4~; \omterm{def:m1:the-buy-side:fees}{commission de performance} 12.4~; total 35.9~millions de dollars.
\textbf{7.} 145.2~millions de dollars~; 3.8~\% par an.
\textbf{8.} 192.8~millions de dollars~; 6.8~\% par an.
\textbf{9.} $92.8 - 35.9 - 45.2 = 11.7$~: la croissance que les frais auraient
rapportée s'ils étaient restés investis. Les frais sont versés tôt et se composent contre
l'investisseur.
\textbf{10.} Le gérant a perçu 39~\% du profit brut~; l'investisseur
a conservé 49~\%~; 12~\% représentent l'effet de composition perdu.
\textbf{11.} $100 \times 1.0595^{10} = \$178.2$ millions, contre 145.2
millions de dollars.
\textbf{12.} Diversification~: un flux non corrélé à ses actions réduit
la volatilité de l'ensemble de la dotation, et il a moins perdu lors des années 2 et 8 si
c'étaient de mauvaises années boursières. Savoir si cela vaut 33 points de pourcentage sur
une décennie est la question du conseil.
\textbf{13.} Rendement brut moyen 8~\%~; ratio de Sharpe $0.08/0.17 = 0.47$~;
$4/0.47^2 = 18$ ans.
\textbf{14.} Valeur liquidative 116.8 contre un \omterm{def:m1:the-buy-side:fees}{high-water mark} de 149.7~: une hausse de 28~\%.
\textbf{15.} (a) L'option est très hors de la monnaie et gagne de la valeur avec
la volatilité~; un risque supplémentaire ne coûte rien de plus au gérant s'il échoue, tandis que
les investisseurs supportent la perte. (b) Un nouveau fonds démarre avec un \omterm{def:m1:the-buy-side:fees}{high-water mark} au pair~: le
gérant échappe à l'écart de 28~\%, et les anciens investisseurs perdent les années «~gratuites~»
qu'elles avaient payées de leurs pertes.
\textbf{16.} Son \omterm{def:m1:the-buy-side:fees}{high-water mark} est sa valeur de souscription. En année 9~: \omterm{def:m1:the-buy-side:fees}{frais de gestion} 2.00,
$N' = 113.00$, \omterm{def:m1:the-buy-side:fees}{commission de performance} $0.2 \times 13 = 2.60$ pour 100 --- alors que le
premier investisseur ne paie aucune \omterm{def:m1:the-buy-side:fees}{commission de performance} cette année-là.
\textbf{17.} Les deux investisseurs détiennent le même portefeuille et doivent des
frais différents~; un \omterm{def:m1:the-buy-side:fees}{high-water mark} unique au niveau du fonds ferait payer au premier investisseur
deux fois les mêmes gains, ou laisser le second en profiter gratuitement.
\textbf{18.} Environ 2.4~\% par an avec une volatilité de 20~\%, contre 1.8~\% avec une
volatilité nulle~: environ 0.6~\% par an pour assumer un risque sans aucune compétence.
\textbf{19.} \textbf{39~\%.}
\textbf{20.} Aucun des deux chiffres en premier~: les modalités qui façonnent le comportement --- le
\omterm{def:m1:the-buy-side:fees}{high-water mark} et ses conditions de réinitialisation, un hurdle, le capital propre du gérant
investi dans le fonds et les limites de risque --- puis les \omterm{def:m1:the-buy-side:fees}{frais de gestion}, qui
sont la partie payée pour rien.
\end{solution}

\begin{solution}{iq:m1:the-buy-side:1}
Le \omterm{def:m1:the-buy-side:benchmark}{tracking error} est la volatilité de l'écart entre le rendement du portefeuille
et celui de l'\omterm{def:m1:the-buy-side:benchmark}{indice de référence}. Une position de 15~\% dans un titre qui pèse 2~\% de
l'indice est un poids actif de 13~\%~: avec une volatilité spécifique de 30~\%, elle ajoute à elle seule
environ 4~\% de \omterm{def:m1:the-buy-side:benchmark}{tracking error}, sans doute plus que ce que le \omterm{def:m1:the-buy-side:mandate}{mandat} autorise, et un
mauvais trimestre sur ce titre fait perdre le \omterm{def:m1:the-buy-side:mandate}{mandat}, quelle que soit la justesse
de la vision à long terme. Pour un fonds public, un plafond légal de 10~\% s'applique de toute façon.
\iqlookfor{poids actif et non poids absolu, et risque de carrière comme une
contrainte réelle.}
\end{solution}

\begin{solution}{iq:m1:the-buy-side:2}
$t = \mathrm{SR}\sqrt{T}$, donc $t = 2$ demande 4 ans dans des
conditions idéales. En pratique, il en faut davantage~: les rendements ne sont ni indépendants ni gaussiens, et si
la stratégie a été choisie parmi beaucoup d'autres, le seuil doit être bien plus élevé que 2.
\iqlookfor{la formule immédiatement, puis la mise en garde sur les tests multiples
sans qu'on la sollicite.}
\end{solution}

\begin{solution}{iq:m1:the-buy-side:3}
Non. Le fonds indiciel trade en raison de flux ou d'un changement
d'indice~: aucune information sur le titre, donc la décote mécanique s'applique
(\cref{prop:m1:the-sell-side:discount}). Le \omterm{def:m1:what-a-trading-firm-does:hedge-fund}{hedge fund} peut vendre sur
une conviction ou juste avant une annonce~: ajouter une charge de \omterm{def:m1:what-a-trading-firm-does:adverse-selection}{sélection adverse}
estimée à partir de ce qui s'est passé après ses précédents ordres. La valeur
de la relation peut jouer en sens inverse.
\iqlookfor{tarifer le client, et pas seulement l'ordre.}
\end{solution}

\begin{solution}{iq:m1:the-buy-side:4}
«~On verse au gérant une part des profits, mais seulement sur les nouveaux profits~: si le
fonds baisse, il doit remonter à son ancien plus haut avant que cette part
soit de nouveau versée.~» Pour le gérant, c'est une option d'achat (call) dont le prix
d'exercice est le plus haut~: près de la valorisation, elle se comporte normalement~; loin en dessous, elle
vaut presque rien et profite de la volatilité, si bien que le gérant est tenté de prendre des risques ou de fermer le
fonds et de le rouvrir avec une nouvelle valorisation.
\iqlookfor{une première phrase simple, puis l'argument de l'option.}
\end{solution}

\begin{solution}{iq:m1:the-buy-side:5}
Reconstitutions d'indices~: chaque fonds indiciel doit acheter les ajouts et vendre les retraits
à la clôture effective~; celui qui achète après l'annonce et vend dans
cette clôture encaisse le gain, aux dépens des fonds indiciels. Les rééquilibrages périodiques
et remises à zéro des fonds à effet de levier~: flux de fin de mois et de fin de journée
de signe et de taille connus~; les fournisseurs de liquidité se positionnent en amont. Aussi
les enchères de clôture elles-mêmes, qui concentrent ce volume.
\iqlookfor{des flux prévisibles en signe, en taille \emph{et} en calendrier.}
\end{solution}

\begin{solution}{iq:m1:the-buy-side:6}
Après trois ans, le ratio de Sharpe estimé d'une stratégie sans compétence est d'environ
$N(0, 1/\sqrt3) = N(0, 0.58)$. Le maximum de 1\,000 normales indépendantes est
d'environ 3.2 écarts-types~: un ratio de Sharpe proche de 1.9. Un
ratio de Sharpe de 2 issu d'un backtest, retenu parmi mille essais, n'a rien de crédible.
\iqlookfor{le maximum espéré de $n$ normales, environ $\sqrt{2\ln n}$,
et la conclusion sur les pratiques de recherche.}
\end{solution}
